\begin{lemma}[Undoing a parallel-extension switch preserves stability]
Let the positive slot \(q.\mathrm{succ}\), representing paper round
\(n=q+1\), be a PE with supporting vertex \(v\), stored gamma
\(\gamma\), and stored delta \(\delta\).  Let \(\sigma_n\) be a proper
coloring leaving exactly \(e_0\) uncolored and stable for
\((T_n,D_n,\varphi_n)\).  Suppose that the
\((\gamma,\delta)\)-component of \(v\) under \(\sigma_n\) meets
\(V(T_n)\) exactly in \(v\).  Define
\[
 \sigma_{n-1}=\sigma_n/P_v(\gamma,\delta,\sigma_n).
\]
Then the PE certificate supplies a selected coloring \(\pi_{n-1}\) and
a second stable coloring \(\pi'_{n-1}\) such that
\(\pi_{n-1}\) is \((T_n,D_{n-1},\varphi_{n-1})\)-stable,
\(\pi'_{n-1}\) is
\((T_n,D_{n-1}\cup\{\delta\},\pi_{n-1})\)-stable, and
\[
 \varphi_n=\pi'_{n-1}/P_v(\gamma,\delta,\pi'_{n-1}).             \tag{1}
\]
Moreover, \(\sigma_{n-1}\) is successively stable relative to
\(\pi'_{n-1}\) and to \(\pi_{n-1}\) for the enlarged control set
\(D_{n-1}\cup\{\delta\}\), and it is
\((T_n,D_{n-1},\varphi_{n-1})\)-stable.
All component switches are on edge identities in the incidence graph.
\end{lemma}
\begin{proof}
Write \(X=V(T_n)\).  Take the certificate at the positive slot
\(q.\mathrm{succ}\) from \noderef{TashkinovSeries}.  It cannot be the
initial constructor.  The PE-kind equation rules out the SE and RE
constructors of \noderef{TashkinovSeriesStep}.  In the remaining PE
constructor, injectivity of \(\mathrm{some}\), applied to the displayed
chosen-vertex, gamma, and delta equations, identifies the constructor's
support, PE color, and selected defective color with
\(v,\gamma,\delta\).  Put
\[
 \pi=\pi_{n-1}=\mathsf{choice.selected},\qquad
 \rho=\pi'_{n-1}=\mathsf{secondStable},\qquad
 f_n=\mathsf{choice.edge}.
\]
The fields of \noderef{MaximumDefectiveChoice} say that \(\pi\) is
proper, leaves exactly \(e_0\) uncolored, is
\((T_n,D_{n-1},\varphi_{n-1})\)-stable, and has an individual edge
\(f_n\) incident with \(v\) for which \(\pi(f_n)=\delta\).  Its defective
vertex field also gives \(v\in X\) by \noderef{DefectiveVertex}.  The PE
constructor says that
\(\rho\) is proper, leaves exactly \(e_0\) uncolored, is
\((T_n,D_{n-1}\cup\{\delta\},\pi)\)-stable, and satisfies
\[
 \varphi_n=\rho/P_v(\gamma,\delta,\rho).                         \tag{2}
\]
Thus \(\pi,\rho\) are the required existential witnesses.

We first determine the component switched in (2).  The rule
\noderef{ParallelExtensionColor} gives
\(\gamma\in\overline\pi(v)\).  By \noderef{MissingColors}, this says
that \(\gamma\) is not in the present-color set of \(v\), while
\noderef{PresentColors} and the incidence and color equations for \(f_n\)
put \(\delta\) in that present-color set.  Consequently
\(\gamma\ne\delta\).  The vertex clause of \noderef{StableColoring}
transports the missingness of \(\gamma\) from \(\pi\) to \(\rho\).
Moreover, \(f_n\) is incident with \(T_n\) by
\noderef{TreeIncidentEdge}, and \(\delta\) belongs to the enlarged
control set.  The edge clause of the same stability therefore gives
\[
 \rho(f_n)=\pi(f_n)=\delta.                                     \tag{3}
\]

In the incidence graph \noderef{TwoColorIncidenceGraph}, every neighbor
of the vertex-node \(v\) is the identity of an edge incident with \(v\)
and colored \(\gamma\) or \(\delta\).  There is no incident
\(\gamma\)-edge because \(\gamma\) is missing at \(v\) under \(\rho\):
otherwise \noderef{PresentColors} would put \(\gamma\) in the complement
of \noderef{MissingColors}.  If an incident \(\delta\)-edge \(a\) were
different from \(f_n\), then \(a,f_n\) would be adjacent by
\noderef{AdjacentEdges} and would have the same color, contrary to the
properness condition \noderef{IsProperEdgeColoring} for \(\rho\).
Equation (3) supplies the one remaining neighbor.  Hence \(v\) has degree
exactly one in this two-color incidence graph, and its unique incident
edge-node is \(f_n\).

The PE constructor supplies an exit witness \noderef{TreeExit}; let \(P\)
be its path from an outside vertex \(u\) to \(v\).  By
\noderef{TreeExitPath}, \(P\) is simple, has positive multigraph-edge
length, has no tree vertex-node other than \(v\), and at least one of
\(\gamma,\delta\) is missing at \(u\).  The first edge-node of \(P\) is
incident with \(u\) and has one of these two colors by
\noderef{TwoColorIncidenceGraph}.  It cannot have a color missing at
\(u\), by the complementary definitions \noderef{MissingColors} and
\noderef{PresentColors}; it therefore has the other color.  Properness
and \noderef{AdjacentEdges} allow at most one incident edge of that other
color, while missingness allows none of the missing color.  Thus \(u\),
like \(v\), has incidence-graph degree exactly one.

Every incidence-graph node has degree at most two by
\noderef{TwoColorIncidenceDegreeAtMostTwo}.  The path \(P\) exhausts the
full component of \(v\): otherwise take a shortest connection from a node
outside \(P\) to \(P\), and consider its first node on \(P\).  At an
internal node, its two distinct path neighbors plus the outside predecessor
would give three neighbors; at either endpoint, the path neighbor plus the
outside predecessor would give degree at least two, contradicting the
degree-one calculation.  Thus the support \(K\) of
\noderef{TwoColorComponent} through \(v\) is exactly the support of \(P\).

We now show that \(f_n\) is the only edge-node of \(K\) incident with
\(T_n\).  Let \(a\in K\) be such an edge-node.  By
\noderef{TreeIncidentEdge}, choose \(w\in X\) incident with \(a\).  The
description in \noderef{TwoColorComponent} makes \(a\) a
\(\gamma\)- or \(\delta\)-colored edge, so the incidence rule in
\noderef{TwoColorIncidenceGraph} makes the vertex-node \(w\) adjacent to
\(a\) and hence a member of \(K=P\).  The exact tree-vertex intersection
clause of \noderef{TreeExitPath} gives \(w=v\).  The uniqueness at \(v\)
proved above then gives \(a=f_n\).  Consequently (2), unfolded using
\noderef{TwoColorComponentSwap}, changes only \(f_n\) among all edges
incident with \(T_n\), changes its value from \(\delta\) to \(\gamma\),
and leaves every other such edge fixed:
\[
 \varphi_n(f_n)=\gamma,\qquad
 \rho(a)=\varphi_n(a)\quad(a\ne f_n\text{ incident with }T_n).  \tag{4}
\]
It also makes \(\delta\) missing at \(v\) under \(\varphi_n\).  Indeed,
the switched edge \(f_n\) now has color \(\gamma\); every other edge at
\(v\) lies outside \(K\), is unchanged from \(\rho\), and cannot have
color \(\delta\) by properness together with (3).  Thus no edge incident
with \(v\) has color \(\delta\), and
\noderef{PresentColors} and \noderef{MissingColors} give
\[
 \delta\in\overline\varphi_n(v).                                \tag{5}
\]

Before analyzing the component under \(\sigma_n\), we verify that the
edge \(f_n\) is controlled by the assumed stability.  At the positive PE
slot, \noderef{TashkinovSeriesStep} gives the exact identity
\(S_n=\{\delta,\gamma\}\).  In the definition
\noderef{TashkinovStoredColors}, \(D_n\) is the union of all \(S_h\) for
\(h\le n\), with \(\overline\varphi_n(T_n)\) removed.  Hence
\(\gamma\) belongs to that cumulative union.  If it is tree-missing, it
belongs to \(\overline\varphi_n(T_n)\); if it is not tree-missing, the
set-difference definition puts it in \(D_n\).  Therefore
\[
 \gamma\in\overline\varphi_n(T_n)\cup D_n.                      \tag{6}
\]
The incidence of \(f_n\) with \(v\in X\) makes it tree-incident by
\noderef{TreeIncidentEdge}.  Equations (4), (6), and the edge clause of
the assumed \noderef{StableColoring} now yield
\[
 \sigma_n(f_n)=\varphi_n(f_n)=\gamma.                            \tag{7}
\]
The vertex clause of that stability transports (5) to
\(\delta\in\overline\sigma_n(v)\).

Let \(L\) be the \((\gamma,\delta)\)-component of \(v\) under
\(\sigma_n\).  Equation (7), the incidence of \(f_n\) with \(v\), and
\noderef{TwoColorIncidenceGraph} put the edge-node \(f_n\) in \(L\).
Suppose an edge-node \(a\in L\) is incident with \(T_n\).  Choose a tree
endpoint \(w\in X\) using \noderef{TreeIncidentEdge}.  Since an edge-node
in the component of the distinct vertex-node \(v\) is connected by a
nonempty incidence walk, \noderef{TwoColorIncidenceGraph} makes \(a\) a
\(\gamma\)- or \(\delta\)-colored edge and makes \(w\) adjacent to \(a\).
Thus \(w\in L\), and the exact intersection hypothesis in the lemma gives
\(w=v\).  At \(v\), no edge has color \(\delta\) because that color is
missing there under \(\sigma_n\); and no edge distinct from \(f_n\) has
color \(\gamma\), since such an edge would be adjacent to \(f_n\) by
\noderef{AdjacentEdges} and violate the properness of \(\sigma_n\) in
\noderef{IsProperEdgeColoring}.  Hence \(a=f_n\).  This proves that
\(f_n\) is also the unique component edge incident with the tree for the
\(\sigma_n\)-component.

Put
\[
 C=D_{n-1}\cup\{\delta\},\qquad
 \tau=\sigma_n/P_v(\gamma,\delta,\sigma_n).
\]
The last uniqueness result and \noderef{TwoColorComponentSwap} show that
the undo switch changes only \(f_n\) among tree-incident edges, changing
its value from \(\gamma\) to \(\delta\).  Together with (4), this gives
\[
 \tau(a)=\sigma_n(a),\quad \rho(a)=\varphi_n(a)
       \quad(a\ne f_n\text{ incident with }T_n),                \tag{8}
\]
and \(\tau(f_n)=\rho(f_n)=\delta\).

We next establish the control-set inclusion needed to compare (8).  At
each \(x\in X\), the vertex clauses of the certificate stabilities give
\[
 \overline\rho(x)=\overline\pi(x)
                  =\overline\varphi_{n-1}(x).                  \tag{9}
\]
Unfolding \noderef{TreeMissingColors} as the union of these vertexwise
missing sets over \(x\in X\) gives
\(\overline\rho(T_n)=\overline\varphi_{n-1}(T_n)\).  Unfolding the same
definition at (5), and using \(v\in X\), gives
\(\delta\in\overline\varphi_n(T_n)\).  Now let a color \(c\) belong to
\(\overline\rho(T_n)\cup C\).  If
\(c\in\overline\rho(T_n)\cup D_{n-1}\), equation (9) and the precise
first inclusion of \noderef{HkuMissingStoredMonotone} put \(c\) in
\(\overline\varphi_n(T_n)\cup D_n\).  The only remaining case is
\(c=\delta\), which was just put in
\(\overline\varphi_n(T_n)\).  Thus
\[
 \overline\rho(T_n)\cup C
 \subseteq \overline\varphi_n(T_n)\cup D_n.                    \tag{10}
\]

We record all coloring hypotheses before using the stability criteria.
The PE constructor in \noderef{TashkinovSeriesStep} gives properness and
\noderef{IsOneEdgeUncolored} for both
\(\varphi_{n-1}\) and \(\varphi_n\);
\noderef{MaximumDefectiveChoice} gives them for \(\pi\); the PE
second-stable fields give them for \(\rho\); and they are hypotheses of
the lemma for \(\sigma_n\).  Since \(\gamma\ne\delta\),
\noderef{TwoColorComponentSwapProper} applied to the proper
\(\sigma_n\) makes \(\tau\) proper.  Finally,
\noderef{TwoColorComponentSwap} leaves an uncolored value uncolored and
maps every colored value to a colored value.  Expanding
\noderef{IsOneEdgeUncolored}, the equation
\(\sigma_n(e_0)=\mathsf{none}\) is therefore preserved, and every
edge different from \(e_0\) remains colored.  Thus \(\tau\) also leaves
exactly \(e_0\) uncolored.

Apply the symmetric criterion \noderef{sc1} first to the proper,
exactly-one-uncolored pair \((\varphi_n,\sigma_n)\) and its assumed
stability.  Let \(a\ne f_n\) be tree-incident.  If either the
\(\rho\)-value or the \(\tau\)-value of \(a\) is controlled by
\(\overline\rho(T_n)\cup C\), equations (8) and inclusion (10) turn this
into the corresponding symmetric controlled-value condition for
\((\varphi_n,\sigma_n)\).  Hence
\(\sigma_n(a)=\varphi_n(a)\), and (8) gives
\(\tau(a)=\rho(a)\).  At \(a=f_n\), both values are \(\delta\).
The same criterion \noderef{sc1}, now applied to the proper,
exactly-one-uncolored pair \((\rho,\tau)\), proves
\[
 \tau\text{ is }(T_n,C,\rho)\text{-stable}.                     \tag{11}
\]

The PE certificate already gives that \(\rho\) is
\((T_n,C,\pi)\)-stable.  Transitivity in \noderef{sc2} applies to the
proper, exactly-one-uncolored triple \((\pi,\rho,\tau)\); together with
(11), it proves that \(\tau\) is \((T_n,C,\pi)\)-stable.  Removing
\(\delta\) from the control set only weakens the edge condition and does
not change the vertex condition in \noderef{StableColoring}, so \(\tau\)
is \((T_n,D_{n-1},\pi)\)-stable.  Finally the selected-coloring field
makes \(\pi\) \((T_n,D_{n-1},\varphi_{n-1})\)-stable.  Transitivity in
\noderef{sc2}, this time for the proper, exactly-one-uncolored triple
\((\varphi_{n-1},\pi,\tau)\), gives
\[
 \tau\text{ is }(T_n,D_{n-1},\varphi_{n-1})\text{-stable}.
\]
Equations (2), (11), and the two transitivity conclusions are exactly the
three asserted stability conclusions for the recovered witnesses.
\end{proof}
